% TOP_PROBLEM: {"id": 2938, "title": "Kirby Problem 4.62", "category_id": 7, "category": {"id": 7, "name": "topology", "display_name": "Topology", "description": "Properties preserved under continuous deformations.", "slug": "topology", "order_index": 7, "created_at": "2026-07-31T15:26:25.671Z"}, "status": "open"}
% FIRST_POSTED: null
% ENTRY_KIND: statement_only
% Imported from: batch13.tex; UnsolvedMath ID 2938
\documentclass[11pt]{article}
\usepackage[margin=25mm]{geometry}
\usepackage{fontspec}
\setmainfont{Latin Modern Roman}
\usepackage{amsmath,amssymb,mathtools}
\usepackage{unicode-math}
\setmathfont{Latin Modern Math}
\usepackage{xurl,hyperref}
\hypersetup{colorlinks=true,urlcolor=blue}
\setcounter{secnumdepth}{0}
\setlength{\parindent}{0pt}
\setlength{\parskip}{5pt}
\setlength{\emergencystretch}{3em}
\newcommand{\sref}[2]{\hyperlink{source:#1:#2}{[S#2]}}
\newcommand{\cataloguescope}[1]{\paragraph{Recorded target.}#1}
\begin{document}
% UnsolvedMath ID 2938; source catalogue code KP-4.62
\setcounter{section}{2937}
\section{Kirby Problem 4.62}\label{2938}
{\small\sffamily UnsolvedMath ID 2938 · Topology}
\subsection{Definitions and mathematical statement}

\paragraph{Shared notation.}$\mathbb N=\{1,2,\ldots\}$, and $\mathbb Z,\mathbb Q,\mathbb R,\mathbb C$ denote the integers, rationals, reals and complex numbers. Any other conventions are specified locally.


Let $X$ be a closed connected oriented smooth four-manifold with $b_2^+(X)>1$, and choose a homology orientation to fix invariant signs. Its intersection form on $H^2(X;\mathbb R)$ has positive index $b_2^+$. Write $\chi(X)=\sum_j(-1)^j\dim H_j(X;\mathbb Q)$ for its Euler characteristic and $\sigma(X)=b_2^+-b_2^-$ for its signature.

Let $x\in H_0(X;\mathbb Q)$ be the class of a point and let $D_X$ denote the $SU(2)$ Donaldson polynomial invariant, obtained by evaluating products of the classes $\mu(a)$ on compactified moduli spaces of anti-self-dual connections on principal $SU(2)$ bundles. The polynomial inputs lie in the symmetric algebra on $H_0(X;\mathbb Q)\oplus H_2(X;\mathbb Q)$; $\mu(x)$ has cohomological degree four and $\mu(a)$ for $a\in H_2$ has degree two. Donaldson simple type means
\[
D_X(x^2z)=4D_X(z)\qquad\text{for every polynomial input }z.
\]
This is the two-point-insertion condition from the original definition. In the bundle-index notation, increasing $c_2$ by one adds eight to the moduli dimension, so two degree-four point insertions are required. \sref{2938}{3}

A spin$^c$ structure $\mathfrak s$ has a determinant line bundle with first Chern class $c_1(\mathfrak s)\in H^2(X;\mathbb Z)$. It is a Seiberg--Witten basic structure if the ordinary integer-valued invariant $SW_X(\mathfrak s)$ is nonzero. This scalar invariant is the signed count in dimension zero and, in nonnegative even dimension $d$, the evaluation of the $d/2$ power of the degree-two point class on the monopole moduli space; it is zero by convention in negative or odd expected dimension. \sref{2938}{4} Seiberg--Witten simple type means that every basic structure has expected dimension
\[
d_X(\mathfrak s)=\frac{c_1(\mathfrak s)^2-2\chi(X)-3\sigma(X)}4=0.
\]
Here the square is evaluated on $[X]$.
\paragraph{Statement recorded in the supplied source.}
For every $X$ as above, answer both questions: (a) does $X$ have Donaldson simple type, and (b) does $X$ have Seiberg--Witten simple type? Simple connectivity, vanishing first Betti number, and the existence of a symplectic structure are not assumed. \sref{2938}{2}
\subsection{Short English statement}
Must every closed oriented smooth four-manifold with positive intersection index greater than one have both Donaldson and Seiberg–Witten simple type?
\subsection{Sources}
\begin{description}
\item[\hypertarget{source:2938:1}{S1}] ULAM.AI and the UnsolvedMath Contributors, UnsolvedMath: A Curated Collection of Open Mathematics Problems, record 2938 (KP-4.62). CC BY 4.0. \url{https://huggingface.co/datasets/ulamai/UnsolvedMath}.
\item[\hypertarget{source:2938:2}{S2}] R. İnanç Baykur, Robion C. Kirby and Daniel Ruberman, K3—A New Problem List in Low-Dimensional Topology, Mathematical Surveys and Monographs 295 (2026), author version, Problem 4.62 and Remarks, PDF page 240. \url{https://math.berkeley.edu/sites/default/files/surv-295-ruberman-watermarked-author-pdf.pdf}.
\item[\hypertarget{source:2938:3}{S3}] Ronald Fintushel and Ronald J. Stern, The Blowup Formula for Donaldson Invariants, Annals of Mathematics 143 (1996), 529--546; arXiv:alg-geom/9405002v2, Section 1, pages 1--3, and Section 5. \url{https://arxiv.org/pdf/alg-geom/9405002v2}.
\item[\hypertarget{source:2938:4}{S4}] Tsuyoshi Kato, Nobuhiro Nakamura and Kouichi Yasui, The simple type conjecture for mod 2 Seiberg–Witten invariants, arXiv:2009.06791v2 (2021), Introduction, Conjecture 1.1 and its following definition. \url{https://arxiv.org/pdf/2009.06791v2}.
\end{description}
\label{end:2938}
\end{document}
